\section{Discussion}
\label{sec:discussion}
The geometric comparisons establish which dimensions remain inconsistent after the source-power fit. The property comparisons then locate the response within the rear temperature field, where boundary position, solidus--liquidus separation and material passage time describe different aspects of solidification. Their relationship determines both the interpretation of the calculated sensitivities and the measurements needed to assess them physically.

\subsection{Constraints supplied by melt-pool geometry}
Fitting the effective source-power fraction fixes the energy supplied by the adopted Gaussian to reproduce B's thermographic mean length. The resulting wide, shallow fusion envelope shows that this axial constraint leaves the transverse dimensions unresolved. Width and depth add information about the distribution of melting across the track; their discrepancies cannot be removed by interpreting the length fit as agreement of the entire fusion boundary. The fitted fraction is conditional on the Gaussian approximation, thermophysical functions, phase law and temperature-threshold length definition. Relating it to physical laser absorption requires an assessment of those choices and of the radiance-based length operator used by the experiment~\cite{lane2020}.

The closest AMMT study illustrates why these choices matter. With the measured beam profile, its isotropic calculation reproduced length more closely than width and depth; calibrated directional conductivity subsequently improved the dimensions~\cite{kollmannsberger2019}. Its isotropic shape was narrower and deeper than the reported metallographic means, whereas the present calculation is wider and shallower. The two calculations also differ in beam representation, property data and phase law. The predecessor identified the limits of fitting three dimensions with a single absorptivity parameter. Here, the transverse residuals identify the shape left unresolved by the B-length fit. The opposite residual directions prevent assigning the present discrepancy to a mechanism inferred from the predecessor. Within the present four-case B comparison, retaining the same source and phase law instead gives the property-extrapolation differences a defined model-specific meaning.

Geometric constraints can also leave thermal predictions unresolved even when more source parameters are fitted. Myers et al. obtained similar cross-sectional area and width-to-depth ratio agreement for several Fresnel/accommodation parameter combinations in bare 316L calculations, but two-color temperature profiles distinguished those combinations~\cite{myers2023}. Their study tests alternative fitted parameter sets. Here, the source fraction is retained after the baseline length fit, so the four IN625 cases measure the response to prescribed property functions. These are complementary designs: the former tests what geometric fitting leaves ambiguous, and the latter identifies how specified material assumptions change the fitted field.

\subsection{Effects of property extrapolation on rear phase boundaries}
Replacing linear extrapolation by constant endpoint values produces opposing responses of conductivity and sensible specific heat. Their combined length change is dominated by the conductivity contribution, whose rear solidus and liquidus displacements are almost equal. Resolving both positions shows that the longer pool principally reflects a downstream shift of the rear temperature profile. The much smaller change in their separation describes a different response: the spatial extent over which the adopted liquid fraction falls from one to zero. The persistence of the width excess and depth deficit across LL, HL, LH and HH further separates this axial response from correction of the transverse fusion shape.

The relative sizes of the imposed changes help interpret the unequal responses. Conservative enthalpy separates sensible storage from the adopted latent contribution~\cite{vanelsen2007}. At 1320~$^\circ$C, constant endpoint extrapolation reduces conductivity by 22.48\% and sensible specific heat by 7.09\%. Because the phase law contributes 4666.67~\si{\joule\per\kilogram\per\kelvin} to effective specific heat in this interval, effective specific heat decreases by only 0.95\%. Integrating the modified specific-heat function from the cold reference also reduces enthalpy at this temperature by 0.66\%. The conductivity and storage cases therefore apply unequal perturbations to diffusive transport and total thermal storage. Their effect sizes characterize these particular functions and the retained phase law. The specific-heat intervention also begins below the solidus, changing pre-melting enthalpy and the temperature--enthalpy inverse along with sensible storage in the freezing interval.

The conductivity response also involves the temperature gradients and the region over which conduction is evaluated. In the HL calculation, the median axial-face P\'eclet diagnostic increases and outward conductive power from the rear cells selected by $T\geq T_s$ and $\xi<0$ rises from 14.724 to 14.913~W. Lower conductivity thus accompanies greater integrated outward power over this solution-dependent region. Both its boundary and its gradients change, so that integral cannot isolate reduced heat escape as the cause of the longer pool. In addition, consistent extrapolation changes the conductivity primitive and the recovered surface temperature used to locate the crossings. Comparing conductive power over an identical spatial region would remove the contribution from changing region selection. Separately evaluating surface-temperature recovery would quantify its contribution to the boundary displacement. Together, these diagnostics would narrow the interpretation of the complete calculation's paired-boundary response.

\subsection{Solidification interval in material coordinates}
Rear-boundary displacement and separation acquire different time meanings in a steadily translating field. A common downstream displacement lengthens the pool without proportionally enlarging the interval between its rear liquidus and solidus. At a fixed speed, nearly unchanged separation consequently gives nearly unchanged passage time through the adopted freezing interval. Between B and C, the speed changes as well as the field: their separations differ by only 0.53~\um{}, but the faster C scan gives a 32.90\% shorter passage time. The 49.02\% increase in interval-mean cooling rate follows from the same separation-and-speed relation. These descriptors distinguish the location of solidification relative to the laser from its duration for stationary material traversing the calculated field.

This distinction is consequential when comparing spatial thermography with a temperature history. Hooper measured centerline profiles and fixed-location histories that resolved heating pulses, cooling and turn events~\cite{hooper2018}. Lane et al. used a constant-speed position-to-time conversion for a below-solidus cooling interval~\cite{lane2020}. The AMMT analysis uses a temperature range and radiance operator that differ from the present liquidus-to-solidus calculation. A direct thermal comparison would therefore require a measured trajectory aligned with the adopted interval, scan position and observation response. The existing passage times describe each condition's own modeled field.

Melt motion changes the trajectories relevant to a microstructure calculation. Fluid-tracer calculations for a Ni--20Cr dispersion-strengthened alloy produced narrower near-surface temperature peaks and hotter subsurface excursions than the corresponding conduction histories~\cite{Hou2024Dissolution}. Transient conduction calculations have also related solidification conditions to grain-structure trends in AlSi10Mg laser and IN718 electron-beam cases, with experimental grain-morphology comparison for the IN718 case~\cite{Plotkowski2017Rapid}. These studies establish the importance of the thermal path and interface motion used in a microstructure model. The present straight-track, quasi-steady mapping supplies a stationary-material centerline interval; applying it to liquid-parcel histories, overlapping scans or grain evolution would require the corresponding transport and transient descriptions.

\subsection{Material data and thermal measurements for model assessment}
The numerical extrapolations make the sensitivity explicit, while their physical selection remains a material-data question. The supplier tables stop below the adopted melting interval, and neither defines a measured liquid-state plateau~\cite{specialmetals625}. Composition-based nickel-superalloy estimates include IN625 conductivity and identify limited underlying liquid-state data~\cite{mills2006}. Such estimates or composition- and state-matched measurements provide candidates for replacing the two prescribed extrapolation rules with physically informed functions. Their comparison should retain an explicit phase law: the dynamic-source IN625 study, for example, uses a 1410~K eutectic endpoint and phase-dependent properties that differ from the present supplier-listed interval~\cite{dynamic2024}.

Source deposition remains a separate candidate for the unresolved shape response. The dynamic volumetric model adjusts deposition depth and effective absorption with the calculated pool, and was calibrated on bare IN625 tracks across several beam diameters before application to a layer case~\cite{dynamic2024}. Recent conduction-model assessment also uses the absorption required to reproduce penetration as an attainability diagnostic under specified transport assumptions~\cite{hong2026}. Aligning source input, phase law, material functions and observation definitions is necessary before differences between these approaches can be assigned to source redistribution or transport. Combined thermography and metallography further show that geometry correlations can coexist with systematic errors outside stable conduction melting~\cite{soares2026}; the operating regime belongs in that comparison.

The present extreme temperatures make this scope particularly consequential. LL, HL, LH and HH positive-$y$ surface-sample peaks reach approximately 2697, 3507, 2756 and 3721~$^\circ$C under a balance that omits evaporation, melt flow and surface deformation. These values expose sensitivity to the assumed heat balance and extrapolations; independent thermal measurements are needed to assess their physical credibility. The large 34--43~\um{} length responses exceed the inspected baseline axial sensitivity, whereas submicrometre interactions retain a need for branch-specific three-dimensional refinement. Analytical verification and residual closure support numerical execution within the inspected scope; physical assessment additionally depends on material data and the distinct experimental populations described in Materials and methods.

The resulting assessment has specific observable targets. Transverse metallography tests the fusion shape beyond the fitted length. Rear thermal profiles test the locations and separation of the adopted phase boundaries. A co-registered temperature history spanning that interval tests the derived material time. Keeping these roles distinct connects the constitutive sensitivity to evidence that can strengthen the calibrated model, rather than treating every geometric or derived thermal descriptor as an interchangeable validation target.

\section{Conclusions}
High-temperature property extrapolation is a consequential assumption in length-calibrated IN625 heat-conduction modeling because the available tables stop below the adopted melting interval. The present comparison makes its consequences explicit for fusion geometry, rear phase-boundary position and the passage time through that interval. Fitting one effective source-power fraction to B thermographic length leaves the calculated fusion envelope wider and shallower than the NIST metallographic means. The other dimensions therefore remain necessary constraints on the spatial distribution of melting.

At retained source power, constant endpoint extrapolation of conductivity lengthens the pool, while the corresponding specific-heat extrapolation shortens it; the combined increase is 34.28~\um{}. Almost equal rear solidus and liquidus displacements account for the dominant conductivity response. Total length primarily changes through movement of this boundary pair, with much less change in its separation and with the transverse discrepancy retained. This distinguishes sensitivity of the rear phase-boundary location from an enlarged solidification interval or an improved fusion shape.

The interval's time interpretation also requires scan speed. Nearly equal B/C rear separations at 0.8 and 1.2~m/s correspond to a roughly one-third shorter passage time and a 49.02\% higher interval-mean cooling rate in C. Together, the results separate the spatial extent of melting, the location and separation of the rear phase boundaries, and the duration of cooling through the adopted interval. Material-specific high-temperature properties and aligned thermal observations can assess these distinct outputs of the calibrated field. The present study supplies their model-based comparison with one retained calibrated source-power fraction and a fixed phase law.
